\subsection{Local probabilities and the conditional GP global studies}
\label{sec:v5-global-results}

The following plots apply the method in Section~\ref{sec:v5-global} to each
campaign's complete mass scan. They distinguish the raw local asymptotic
reference, the local probability after accounting for the archived
background response, and the scan-global probability under that same model.
The global correction acts on a specified probability ordering; it does not
rescale the observed upper limit.

\fig{0.98}{../figures/v5_global_probabilities_2015.pdf}{Profiled 2015 local and
conditional GP global probabilities. The study used ten complete pilot
scans, 1,000 independent validation scans and 200,000 GP fields. The principal
global probability at the selected 51~MeV threshold is 0.018985 in the GP
approximation; 13 of 1,000 direct scans exceed that threshold. The precision
of either ensemble and the generating-background assumptions remain explicit.}
{fig:v5-global-2015}
\fig{0.98}{../figures/v5_global_probabilities_2016.pdf}{Profiled 2016 conditional
global study. The leading principal-ordering threshold at 42~MeV has no
exceedance in either ensemble. The corresponding one-sided 95\% bounds are
$1.498\times10^{-5}$ for 200,000 GP fields and 0.002991 for 1,000 direct scans.
They are finite-sampling bounds under the specified models, not a resolved
rare particle probability. A qualified final 2016 global result remains
outstanding.}
{fig:v5-global-2016}
\fig{0.98}{../figures/v5_global_probabilities_2021.pdf}{Profiled global study for
the released 2021 10\% sample. The principal ordering selects 92~MeV and gives
GP global probability 0.02519, with 26 of 1,000 independent direct scans above
the threshold. This selection differs from the largest raw positive fitted
signal. Qualification of the generating background and the final analysis
choices remains required for a final particle-global interpretation.}
{fig:v5-global-2021}
\FloatBarrier

The combined study illustrates why the distinction matters. At 76~MeV the
observed raw local coordinate is only 0.166, corresponding to asymptotic
$p_0=0.434$. The archived stress spectrum has an offset of $-8.700$ and a
response width of 0.979, giving an offset-adjusted score of about 9.05.
This is a large disagreement with that reference response, not a $9\sigma$
observed resonance. Under the separate raw-maximum ordering, every simulated
scan exceeds the observed maximum. The two tests cannot be interchanged
according to which yields the more striking probability.

\fig{0.98}{../figures/v5_global_probabilities_combined.pdf}{Full combined
search with the same profiled local and conditional global probability
conventions. The upper panel shows the observed local coordinate and the
stress-background offset so that an extreme centered tail can be traced to
its reference. Open downward triangles are one-sided 95\% Monte Carlo upper bounds
for zero counts; triangles at the local plotting floor mark analytic values
below $10^{-8}$. The raw and centered orderings are separate tests.}
{fig:v5-global-combined}
\FloatBarrier

\FloatBarrier\clearpage
\begin{landscape}
\subsection{Agreement of fitted rates across campaigns}
\label{sec:v5-rate-consistency}

A shared coupling requires the signal yields in different campaigns to agree
after their different normalizations are applied. Figure~\ref{fig:v5-consistency}
compares the separate fitted coupling estimates with the common fit. The
quantity formerly labelled $\Delta D$ is the loss in twice the log likelihood
when separate campaign amplitudes are replaced by one shared amplitude:
\begin{equation}
 \Delta D=2\left[\mathrm{NLL}_{\rm common}-
                    \sum_d\mathrm{NLL}_{d,\rm independent}\right].
\end{equation}
A larger loss means the common rate describes the selected local spectra less
well than separate rates. The number of freed amplitudes is one fewer than
the number of active campaigns. Because the displayed masses were selected
from the data, this descriptive comparison is not assigned an unadjusted
compatibility probability.

\methodfig{0.98}{../figures/v5_dataset_amplitude_consistency.pdf}{Separate campaign
coupling estimates and local curvature errors compared with the common rate.
The likelihood loss quantifies how much the local fit improves when each year
can choose its own amplitude. The 66~MeV fit has positive contributions from
all three years; the 92~MeV separate rates differ more visibly. These selected
comparisons and local errors are not post-selection compatibility tests.}
{fig:v5-consistency}
\FloatBarrier

\end{landscape}
\clearpage
